\section{Target, information, and what has to be estimated}\label{sec:target}

\subsection{The two objects}

Write $\theta_c$ for population $c$'s latent summary and $Y_c$ for its survey
estimate, taken as a deviation from a fixed centre. The observed and the latent
object are not interchangeable, and the paper keeps them apart throughout:
\begin{equation}\label{eq:decomp}
Y_c=\theta_c+e_c,\qquad \operatorname{Var}(\theta_c)=A,\qquad
\operatorname{Var}(e_c)=D_c ,
\end{equation}
with $D_c$ the design variance of the survey estimate. The \emph{design share}
\begin{equation}\label{eq:rho}
\rho_c^2=\frac{D_c}{A+D_c}\in[0,1)
\end{equation}
is the fraction of the observed dispersion that is sampling error. An interval
calibrated on the $Y_c$ covers a further \emph{survey estimate}; covering a
further \emph{latent summary} requires removing that fraction.

\subsection{The interval}

Fix an integer $m\le K$, write $\operatorname{ord}_m$ for the $m$th smallest of
$K$ numbers, and let $t=D/A$ be the scale ratio, taken common across populations
for this section. Standardising $Z_c(t)=Y_c/\sqrt{1+t}$, the interval is
\begin{equation}\label{eq:band}
B(t)=[-R(t),\,R(t)],\qquad
R(t)=\operatorname{ord}_m\{|Z_c(t)|\}_{c=1}^K=q\,\kappa,
\end{equation}
where $q=\operatorname{ord}_m\{|Y_c|\}$ and
$\kappa=(1+t)^{-1/2}=\sqrt{A/(A+D)}=\sqrt{1-\rho^2}$. Here $q$ is a
statistic---the radius the observed spread itself supplies---and $\kappa$ is the
factor carrying it to the latent scale.

\begin{lemma}[Base validity]\label{lem:rank}
If $\theta_{K+1},Z_1(t),\dots,Z_K(t)$ are exchangeable with a continuous
distribution, then $\Pr\{\theta_{K+1}\in B(t)\}=m/(K+1)$. Under
\eqref{eq:decomp} with equal design variances this holds exactly at the true $t$,
because $Z_c(t)$ and $\theta_{K+1}$ are then i.i.d.\ $N(0,A)$. Write
$\alpha_0=1-m/(K+1)$.
\end{lemma}

Everything that follows concerns one difficulty: $t$ is unknown, so the interval
is evaluated at a bound for it rather than at its value.

\subsection{What has to be known precisely, and it is not $A$}

\textbf{The band depends on the two scales only through their ratio $\kappa$},
and not on $A$ separately. Since $\widehat A$ and $\widehat{A+D}$ are built from
the same observed spread, their errors partly cancel, and to first order
\begin{equation}\label{eq:elasticity_kappa}
\frac{\operatorname{RSE}(\widehat\kappa)}{\operatorname{RSE}(\widehat{\sqrt A})}
=\sqrt{\frac{\rho^4\sigma_T^2+\sigma_D^2}{\sigma_T^2+\sigma_D^2}}
\ \longrightarrow\ \rho^2
\end{equation}
as the design variances become well determined, where $\sigma_T^2$ and
$\sigma_D^2$ are the variances of the estimated total and design scales. A
requirement stated on $\widehat A$ therefore overstates what the band needs by a
factor $\rho^{-4}$ in the number of populations. Simulation puts the realised
gain at 14 to 24 times at $\rho^2=0.2$ and 2.0 to 3.8 times at $\rho^2=0.5$,
against the predicted $\rho^{-4}$ of 25 and 4.

\subsection{Two axes, not one}

Writing $\sigma_D^2=2D^2/\nu_{\mathrm{tot}}$ for a design-variance estimate on
$\nu_{\mathrm{tot}}$ degrees of freedom, the requirement
$\operatorname{RSE}(\widehat\kappa)\le\varepsilon$ becomes
\begin{equation}\label{eq:twoaxis}
\boxed{\ \frac{1}{K-1}+\frac{1}{\nu_{\mathrm{tot}}}\ \le\
\frac{2\varepsilon^2(1-\rho^2)^2}{\rho^4}\ }
\end{equation}
Populations and design degrees of freedom enter \emph{harmonically}: neither
substitutes for the other, and $\nu_{\mathrm{tot}}$ imposes a floor no number of
populations clears. Replicate count does not help---with design degrees of
freedom $\nu^{\mathrm{des}}=20$, raising a resampling count from 100 to 10{,}000
buys 3.6 effective degrees of freedom and no more.

\paragraph{The boundary is not an artefact of this estimator.} Condition
\eqref{eq:twoaxis} is sufficient, and it was derived from one estimator's
relative standard error. A matching necessary condition holds. Restrict to equal
design variances and take two parameter points separated so that their
$\varepsilon$-balls in $\kappa$ are disjoint. The scaled-$\chi^2$ and normal
Hellinger affinities are available in closed form, so Le Cam's two-point bound \citep{lecam1986}
gives, exactly and in finite samples, that \emph{no} procedure can certify
$\kappa$ to tolerance $\varepsilon$ with failure probability $\eta<1/2$ unless
\begin{equation}\label{eq:necessary}
\nu_{\mathrm{tot}}\ \ge\ L(\eta)\,\frac{\rho^4}{\varepsilon^2(1-\rho^2)^2},
\qquad L(\eta)=-\log\big\{2\sqrt{\eta(1-\eta)}\big\},
\end{equation}
and the same bound with $K$ in place of $\nu_{\mathrm{tot}}$. The dependence on
$\rho$ and on $\varepsilon$ is the same as in \eqref{eq:twoaxis}; the constants
differ by a factor $2.4$ to $2.9$ over the range we checked ($\rho^2\ge0.4$,
$\varepsilon\le0.05$, $\eta=0.05$), one side being exact and the other a normal
approximation. At small design share the perturbation approaches the edge of the
parameter space and the bound loosens, to a factor $14$ at $\rho^2=0.2$,
$\varepsilon=0.05$.

The two axes carry the \emph{same} constant, and not by coincidence. Writing
$\tau$ for the ratio of the two $\kappa$ values, the $\nu$-axis perturbation
multiplies $D$ by $(1-\tau^2\kappa^2)/(1-\kappa^2)$ while the $K$-axis
perturbation multiplies $A+D$ by exactly its reciprocal, and the affinity
$\{2\sqrt r/(1+r)\}^{n/2}$ is invariant under $r\mapsto 1/r$. That is why the two
enter harmonically rather than in some other combination.

\emph{Neither the technique nor, in this submodel, the problem is new.} With
equal design variances the model is the balanced one-way random-effects model,
$\kappa$ is a monotone function of the intraclass correlation, exact intervals
for that coefficient are classical, and trading the number of groups against the
number of observations per group to control interval length is an established
design literature \citep{donner1986,zou2012icc,mondal2024review}. The
minimax theory for signal-to-noise ratios in high-dimensional regression
\citep{verzelen2018snr} concerns a different observation structure: there is no
separate channel carrying the error scale on a known number of degrees of
freedom, which is what produces the second axis here. We state
\eqref{eq:necessary} to establish that \eqref{eq:twoaxis} reflects the problem
rather than our estimator, and we claim nothing further for it.

\subsection{Estimating the design variances costs little; the estimator costs a lot}

Treating $D_1,\dots,D_K$ as unknown nuisance parameters observed through
$\nu_c\Dhat_c/D_c\sim\chi^2_{\nu_c}$, the effective information for $A$ is
\begin{equation}\label{eq:ieff}
I_{\mathrm{eff}}(A)=\tfrac12\sum_c\big\{(A+D_c)^2+D_c^2/\nu_c\big\}^{-1},
\end{equation}
so the standard error for $A$ inflates, relative to known design variances, by at
most $\sqrt{1+\rho^4/\nu_{\min}}$: 3.2\% at $\nu=10$ and $\rho^2=0.8$, 7.7\% at
$\nu=4$, 14.9\% at $\nu=2$. \textbf{At the degrees of freedom we examine, not
knowing the design variances is a second-order cost.} The bound is attained
exactly, in finite samples, by a moment estimator with weights
$w_c\propto\{(A+D_c)^2+D_c^2/\nu_c\}^{-1}$.

The first-order cost is the estimator. Under dispersed design variances the
unweighted subtraction $\widehat A=K^{-1}\sum_c(Y_c^2-\Dhat_c)$ has relative
efficiency 23\% against \eqref{eq:ieff}, where oracle weighting reaches 99\%. And
substituting $\Dhat_c$ into a known-variance likelihood biases $\widehat A$
\emph{upward} by approximately $4\rho^4/\{\nu(1-\rho^2)\}$ relative to $A$, a
bias of order $1/\nu$ that does not average away as populations accumulate. The
second-order mean-squared-error corrections of \citet{prasad1990} and
\citet{datta2000} address the same estimated-variance effect for an in-sample
predictor, where the target makes the correction additive rather than a
multiplicative distortion of a band.

These are properties of the estimation problem. What matters for the interval is
the next section.
